% Generated by tools/edn_latex.py from reports/edn.md.
% Do not edit: change reports/edn.md and run `make edn`.
\ednentry{
  id = {EDN-63},
  title = {\texorpdfstring{\texttt{metrics.\allowbreak{}json}}{metrics.json} carries the data source and the fitted estimator, at the cost of five more columns},
  date = {2026-10-01},
  milestone = {M3: Quality Assurance},
  topic = {Experiment Tracking, Release Gate},
  participants = {@lukas2510},
  decision = {\texttt{metrics.\allowbreak{}json} gains a \texttt{data\_\allowbreak{}source} leaf, the resolved \texttt{download.\allowbreak{}source}, and an \texttt{estimator} leaf per variant, the model actually fitted; each \texttt{reports/\allowbreak{}metrics/\allowbreak{}<variant>.\allowbreak{}json} carries both as well. That widens the DVC metrics file from 24 leaves to 29, and \texttt{evaluate} declares \texttt{download.\allowbreak{}source} under \texttt{params:} so the recorded value cannot disagree with the lock.},
  rationale = {\emph{Alternatives considered:}
\begin{itemize}[leftmargin=*, topsep=2pt]
\item \textbf{Option A (chosen): the provenance goes into \texttt{metrics.\allowbreak{}json} itself, five leaves wide.}
\newline Pros: \texttt{dvc metrics show} and \texttt{dvc metrics diff} are where a promotion reviewer reads a result (EDN-12), so the one place a number can be quoted from is the one place that now says what produced it. \texttt{dvc metrics diff} lists a changed leaf, so a source or estimator that changes between two commits shows up beside the metric it moved.
\newline Cons: the file was already past the width a terminal table keeps, and this is five more columns of something that is not a measurement.
\item \textbf{Option B: the provenance goes only into the per-variant records under \texttt{reports/\allowbreak{}metrics/\allowbreak{}}.}
\newline Pros: costs nothing in the metrics table; the per-variant record already held \texttt{estimator}.
\newline Cons: it answers the question in the file nobody reads first. The defect this exists to prevent was a reader quoting \texttt{dvc metrics show}, and a second file they have to open is exactly the step that was skipped.
\item \textbf{Option C: an MLflow tag on each run and nothing in the artefacts.}
\newline Pros: the experiment is where the ladder is compared, and a tag is free.
\newline Cons: tracking is optional by construction -- it degrades to disabled with no credentials -- so the one artefact that is always written would be the one that says nothing. It also puts the evidence behind a login, while \texttt{metrics.\allowbreak{}json} is in the pull request diff.
\item \textbf{Option D: nothing, and rely on \texttt{dvc.\allowbreak{}lock} recording the source.}
\newline Pros: the lock already records it, so the information exists.
\newline Cons: it exists in a file whose \texttt{download} stage entry nobody reads to interpret an MdAPE, and it is a different file from the one holding the number. That is the state that let a constant-median stub's figure on generated rows sit on \texttt{main} as \texttt{best\_\allowbreak{}mdape} for a week.
\end{itemize}
\par
\emph{Why:} The honest version of the trade-off is that the metrics file is already too wide to read in a terminal and this makes it wider, so the question is only whether provenance earns five of those columns. It does, because the failure it prevents happened: \texttt{metrics.\allowbreak{}json} on \texttt{main} carried one MdAPE under all four variant names, computed by a constant-median stub on a synthetic fixture, and nothing in the artefact said either thing. The width argument also has a limit that matters here: the file is already read with \texttt{dvc metrics show -\allowbreak{}-\allowbreak{}md}, where five more columns cost nothing a reader notices, while \texttt{dvc metrics diff} is long-format and lists only what changed, so the new leaves are invisible until they are the answer.
Declaring \texttt{download.\allowbreak{}source} under the stage's \texttt{params:} is the part that makes it a guarantee rather than a label. Without it the stage would read a key DVC does not watch, and a recorded value could drift from the lock; with it, the convention \texttt{dvc.\allowbreak{}yaml} states for every other stage holds for this one too.
One consequence was found while implementing it and is worth recording: the end-to-end test fixture ran the stages against the project's own \texttt{params.\allowbreak{}yaml}, so after the flip to \texttt{zenodo} the fixture's artefacts would have claimed \texttt{data\_\allowbreak{}source: zenodo} over generated rows -- the exact mislabel the field exists to prevent. The fixture now runs against a params copy pinned to \texttt{synthetic}, and a test asserts the artefacts say so, which is what keeps the new field from being decorative.},
  ai = {Alternative generation, Alternative assessment, Solution generation},
  response = {Accepted},
  assessment = {Issue \#57 had already identified the defect and asked for the source and the estimator in \texttt{metrics.\allowbreak{}json}, so the problem was not AI's to find. What AI contributed is the shape and the two things that make it hold: declaring \texttt{download.\allowbreak{}source} under the stage's \texttt{params:} rather than reading an undeclared key, and noticing that the test fixture would start mislabelling its own artefacts as soon as the project's source flipped. The second one is the kind of defect this project has repeatedly produced -- a claim the code does not back -- and it would have shipped as a passing test suite.},
  aievidence = {Claude Code session on 2026-10-01 implementing issue \#57: the four options were laid out against the existing docstring's 24-leaf constraint, and the fixture's mislabelling was found by asking what \texttt{params[\textquotedbl{}download\textquotedbl{}][\textquotedbl{}source\textquotedbl{}]} would resolve to inside the test run rather than inside the pipeline.},
  evidence = {\href{https://github.com/mlops-2627q1-mds-upc/MLOPS_Recommenditos/blob/main/recommenditos/modeling/evaluate.py}{\texttt{recommenditos/\allowbreak{}modeling/\allowbreak{}evaluate.\allowbreak{}py}}, \texttt{gate\_\allowbreak{}summary} and \texttt{\_\allowbreak{}evaluate\_\allowbreak{}variant}; \texttt{tests/\allowbreak{}test\_\allowbreak{}pipeline.\allowbreak{}py::test\_\allowbreak{}the\_\allowbreak{}metrics\_\allowbreak{}artefacts\_\allowbreak{}say\_\allowbreak{}which\_\allowbreak{}data\_\allowbreak{}and\_\allowbreak{}which\_\allowbreak{}estimator\_\allowbreak{}produced\_\allowbreak{}them}; \texttt{metrics.\allowbreak{}json}; \hyperref[EDN-12]{EDN-12}, \hyperref[EDN-33]{EDN-33}; \href{https://github.com/mlops-2627q1-mds-upc/MLOPS_Recommenditos/issues/57}{issue \#57}.},
}
